% Randwertproblem: Poisson in V, Dirichlet-Rand (phi vorgegeben) und Neumann-Rand (dphi/dn vorgegeben)
\input{style}
\begin{document}
\begin{tikzpicture}
  \coordinate (P1) at (-2.0,0.3); \coordinate (P2) at (0.2,1.55);
  \coordinate (P3) at (2.2,0.5);  \coordinate (P4) at (0.4,-1.45);
  \fill[soft] (P1) .. controls (-1.9,1.3) and (-0.8,1.6) .. (P2)
    .. controls (1.2,1.5) and (2.2,1.3) .. (P3)
    .. controls (2.2,-0.4) and (1.5,-1.4) .. (P4)
    .. controls (-0.6,-1.5) and (-2.1,-0.8) .. (P1);
  % Dirichlet-Teil (links oben): Akzent
  \draw[acc, line width=2pt] (P4) .. controls (-0.6,-1.5) and (-2.1,-0.8) .. (P1)
    .. controls (-1.9,1.3) and (-0.8,1.6) .. (P2);
  % Neumann-Teil (rechts): blau
  \draw[blue, line width=2pt] (P2) .. controls (1.2,1.5) and (2.2,1.3) .. (P3)
    .. controls (2.2,-0.4) and (1.5,-1.4) .. (P4);
  % Normale auf dem Neumann-Rand
  \draw[vec] (2.03,-0.25) -- ++(0.62,-0.22) node[anchor=west] {$\vv n$};
  % Ladungsverteilung im Inneren
  \fill[warm, fill opacity=0.35] (0.1,0.1) ellipse (0.45 and 0.3);
  \node[warm, font=\footnotesize] at (0.1,0.1) {$\rho$};
  \node[acc] at (-1.1,0.2) {$V$};
  \node[font=\footnotesize] at (0.25,-0.75) {$\nabla^2\phi=-\rho/\varepsilon$};
  \node[acc, anchor=east, align=right, font=\footnotesize] at (-2.15,0.6) {Dirichlet:\\ $\phi|_{\partial V}=\phi_0$};
  \node[blue, anchor=west, align=left, font=\footnotesize] at (2.3,1.15) {Neumann:\\ $\partial\phi/\partial n|_{\partial V}=g$};
\end{tikzpicture}
\end{document}
